% -----------------------------------------------------------------------------
% Electron-photon scattering and loop structures in strict SI units
% -----------------------------------------------------------------------------

真实光子出现在外腿以后，除电子自旋外还必须处理物理偏振、同一电子线上的不同顶角次序以及规范不变性。以下过程都使用\eqnrefs{eq:gem-alpha-si-r10,eq:si-reduced-propagators-r10,eq:si-reduced-vertices-r10} 的 SI 约化规则：\(g_{\rm em}=\sqrt{4\pi\alpha}\) 是无量纲图耦合，四动量保持物理动量量纲。

对电子静止系中的 Compton 散射
\begin{equation*}
 e^-(p)+\gamma(k)\to e^-(p')+\gamma(k'),
\end{equation*}
两个光子顶角在同一电子线上的两种次序必须在振幅层相干相加：
\begin{equation}
 \boxed{
 \mathcal M_C
 =-g_{\rm em}^2\bar u(p')\left[
 \slashed\varepsilon'^{*}
 \frac{\slashed p+\slashed k+m_{\mathrm e} c}{2p\cdot k}
 \slashed\varepsilon
 -
 \slashed\varepsilon
 \frac{\slashed p-\slashed k'+m_{\mathrm e} c}{2p\cdot k'}
 \slashed\varepsilon'^{*}
 \right]u(p).}
 \label{eq:compton-M}
\end{equation}
两图之和满足外部光子的 Ward 条件；单独任一图都不具有这一性质。对未偏振散射，定义两个不变量
\begin{equation*}
 a\equiv p\cdot k,\qquad b\equiv p\cdot k'.
\end{equation*}
完成初态电子自旋和光子偏振平均、末态求和以后，树级振幅给出
\begin{equation}
 \boxed{
 \overline{|\mathcal M_C|^2}
 =2g_{\rm em}^4\left[
 \frac ab+\frac ba
 +2m_{\mathrm e}^2c^2\left(\frac1a-\frac1b\right)
 +m_{\mathrm e}^4c^4\left(\frac1a-\frac1b\right)^2
 \right].}
 \label{eq:compton-M2-unpolarized-dp73}
\end{equation}
这个式子把所有 Dirac 迹和偏振求和压缩成 Lorentz 不变量；于是只需代入具体参考系的运动学与二体相空间。

四动量守恒与电子质量壳给出精确反冲关系
\begin{equation}
 \boxed{
 \frac{\omega'}{\omega}
 =\frac{1}{1+x(1-\cos\theta)},
 \qquad
 x\equiv\frac{\hbar\omega}{m_{\mathrm e} c^2}.}
 \label{eq:compton-energy-ratio}
\end{equation}
等价地
\begin{equation}
 \boxed{
 \frac1{\omega'}-\frac1\omega
 =\frac{\hbar}{m_{\mathrm e} c^2}(1-\cos\theta).}
 \label{eq:compton-shift-SI}
\end{equation}
唯一的反冲控制量是 \(x\)。当 \(x\ll1\) 时回到低能 Thomson 极限；\(x\gtrsim1\) 时必须保留完整反冲运动学。


\begin{exbox}[$17\,\mathrm{keV}$ X 射线 Compton 反冲]
取电子初态静止，入射光子能量 $E_\gamma=\hbar\omega=17\,\mathrm{keV}$。首先计算唯一的反冲控制量
\begin{equation*}
 x=\frac{E_\gamma}{m_{\mathrm e} c^2}
 =\frac{17\,\mathrm{keV}}{510.999\,\mathrm{keV}}
 \simeq 3.33\times10^{-2}.
\end{equation*}
所以这是低能但已经可以分辨反冲的区域。若散射角为 $\theta=90^\circ$，\eqnrefs{eq:compton-energy-ratio} 给
\begin{align*}
 E_\gamma'
 &=\frac{17\,\mathrm{keV}}{1+x}
 \simeq16.45\,\mathrm{keV}.
\end{align*}
用 $\lambda=2\pi c/\omega=h c/E_\gamma$ 把\eqnrefs{eq:compton-shift-SI} 改写成波长形式，得到
\begin{equation}
 \boxed{
 \Delta\lambda
 =\lambda'-\lambda
 =\frac{h}{m_{\mathrm e} c}(1-\cos\theta).}
 \label{eq:compton-wavelength-shift-dp61}
\end{equation}
因此 $90^\circ$ 时
\begin{equation*}
 \Delta\lambda=\frac{h}{m_{\mathrm e} c}=2.426\,\mathrm{pm},
\end{equation*}
而完全后向散射 $\theta=180^\circ$ 时波长移动加倍为 $4.853\,\mathrm{pm}$，光子能量降到约 $15.94\,\mathrm{keV}$。这个例子把“反冲参数 $x$”变成了直接可测的谱线位移：$x\ll1$ 并不意味着位移严格为零，只意味着 Klein--Nishina 角分布可以围绕 Thomson 极限作受控展开。
\end{exbox}

电子静止系的入射通量标量为
\begin{equation}
 \mathcal I_{e\gamma}
 =\sqrt{(p\!\cdot k)^2-m_{\mathrm e}^2c^2k^2}
 =p\!\cdot k
 =m_{\mathrm e}\hbar\omega.
 \label{eq:compton-invariant-flux-r14}
\end{equation}
把振幅、自旋--偏振求和与二体相空间组合，得到 Klein--Nishina 公式
\begin{equation}
 \boxed{
 \frac{\dd\sigma_C}{\dd\Omega}
 =\frac{r_e^2}{2}
 \left(\frac{\omega'}{\omega}\right)^2
 \left(
 \frac\omega{\omega'}+\frac{\omega'}\omega-\sin^2\theta
 \right).}
 \label{eq:klein-nishina-SI}
\end{equation}
其中 \(r_e=\alpha\hbar/(m_{\mathrm e} c)\) 是经典电子半径。详细的相空间雅可比行列式与 Dirac 迹统一放在证明区。

\begin{figure}[htbp]
 \centering
 \begin{tikzpicture}[scale=.76, baseline=(current bounding box.center)]
\begin{scope}[xshift=-4cm]
\coordinate (v1) at (-.7,0);\coordinate (v2) at (.7,0);
\draw[fermion] (-2.2,-1.2)--(v1)--(v2)--(2.2,1.2);
\draw[photon] (-2.2,1.2)--(v1);\draw[photon] (v2)--(2.2,-1.2);
\node[above=3pt,fill=white,inner sep=1pt] at (0,0){$p+k$};
\end{scope}
\node at (0,0){$+$};
\begin{scope}[xshift=4cm]
\coordinate (v1) at (-.7,0);\coordinate (v2) at (.7,0);
\draw[fermion] (-2.2,-1.2)--(v1)--(v2)--(2.2,1.2);
\draw[photon] (-2.2,1.2)--(v2);\draw[photon] (v1)--(2.2,-1.2);
\node[above=3pt,fill=white,inner sep=1pt] at (0,0){$p-k'$};
\end{scope}
\end{tikzpicture}

 \caption{Compton 散射的两张树图。两个光子顶角在同一电子线上有两种次序；只有两图相干相加后才满足 Ward 条件。}
 \label{fig:compton-two}
\end{figure}

\begin{figure}[htbp]
 \centering
 \includegraphics[width=.78\linewidth]{compton_klein_nishina.png}
 \caption{Klein--Nishina 角分布。取 \(x=\hbar\omega/(m_{\mathrm e} c^2)=0.01,0.5,2.0\)；随着 \(x\) 增大，相对论反冲逐渐抑制后向散射。}
 \label{fig:klein-nishina-plot}
\end{figure}

交叉变换把同一局域顶角结构连接到 \(e^-e^+\to\gamma\gamma\)。把 Compton 振幅的一条出射电子外线解析延拓为入射正电子，同时把相应外部四动量反向，得到
\begin{equation}
 \boxed{
 \mathcal M_{\gamma\gamma}
 =-g_{\rm em}^2\bar v(p_2)\left[
 \slashed\varepsilon_2^*
 \frac{\slashed p_1-\slashed k_1+m_{\mathrm e} c}{(p_1-k_1)^2-m_{\mathrm e}^2c^2}
 \slashed\varepsilon_1^*
 +\slashed\varepsilon_1^*
 \frac{\slashed p_1-\slashed k_2+m_{\mathrm e} c}{(p_1-k_2)^2-m_{\mathrm e}^2c^2}
 \slashed\varepsilon_2^*
 \right]u(p_1).}
 \label{eq:ee-gammagamma-M-r68}
\end{equation}
这个式子给出\eqnrefs{eq:compton-M} 所代表的同一解析四点函数在另一组物理质量壳边界值上的表示。
\begin{figure}[htbp]
 \centering
 \begin{tikzpicture}[scale=.76, baseline=(current bounding box.center)]
\begin{scope}[xshift=-4cm]
\coordinate (v1) at (-.7,0);\coordinate (v2) at (.7,0);
\draw[fermion] (-2.2,1.2)--(v1);\draw[antifermion] (-2.2,-1.2)--(v1);\draw[fermion] (v1)--(v2);\draw[photon] (v2)--(2.2,1.2);\draw[photon] (v1)--(2.2,-1.2);
\end{scope}
\node at (0,0){$+$};
\begin{scope}[xshift=4cm]
\coordinate (v1) at (-.7,0);\coordinate (v2) at (.7,0);
\draw[fermion] (-2.2,1.2)--(v1);\draw[antifermion] (-2.2,-1.2)--(v1);\draw[fermion] (v1)--(v2);\draw[photon] (v1)--(2.2,1.2);\draw[photon] (v2)--(2.2,-1.2);
\end{scope}
\end{tikzpicture}

 \caption{\(e^-e^+\to\gamma\gamma\) 的两张电子交换树图。它们与 Compton 两图属于同一解析振幅的不同交叉道。}
 \label{fig:ee-gammagamma}
\end{figure}

树级图中的内部四动量由顶角守恒关系完全固定；圈图则留下不能由外部动量唯一确定的独立圈四动量，需要积分并进行紫外重整化。QED 的三个基本一圈单粒子不可约核是电子自能、光子真空极化与电磁顶角修正；四外光子过程则从电子盒图开始。它们的具体积分、Ward 约束、重整化与形状因子在紧接着的一圈计算中统一建立。



阿贝尔 Maxwell 拉格朗日量没有基本 \(A^3\) 或 \(A^4\) 顶角，因此 \(\gamma\gamma\to\gamma\gamma\) 的最低阶连通振幅来自闭合电子盒图，阶数为 \(g_{\rm em}^4\sim\alpha^2\)。当 \(\hbar\omega/(m_{\mathrm e} c^2)\ll1\) 时，电子圈可以匹配为局域四场强算符；规范不变性要求每条外部光子通过场强张量进入，所以振幅至少含四次外部频率。精确系数由电子盒图的圈积分确定。


\begin{figure}[htbp]
 \centering
 \begin{tikzpicture}[scale=.70, baseline=(current bounding box.center)]
% self-energy
\begin{scope}[xshift=-6.2cm]
\draw[fermion] (-1.8,0)--(1.8,0);
\draw[photon] (-.7,0) .. controls (-.7,1.4) and (.7,1.4) .. (.7,0);
\node[below] at (0,-.45){self-energy};
\end{scope}
% vacuum polarization
\begin{scope}[xshift=-2.0cm]
\draw[photon] (-1.8,0)--(-.8,0);\draw[photon] (.8,0)--(1.8,0);
\draw[fermion] (-.8,0) arc (180:0:.8 and .8);\draw[fermion] (.8,0) arc (0:-180:.8 and .8);
\node[below] at (0,-1.1){vacuum polarization};
\end{scope}
% vertex
\begin{scope}[xshift=2.5cm]
\draw[fermion] (-1.7,-.7)--(0,0)--(1.7,-.7);\draw[photon] (0,0)--(0,1.8);
\draw[photon] (-.55,-.23) .. controls (0,.8) and (.55,.8) .. (.55,-.23);
\node[below] at (0,-1.1){vertex correction};
\end{scope}
% box
\begin{scope}[xshift=6.8cm]
\draw[fermion] (-.7,.7)--(.7,.7)--(.7,-.7)--(-.7,-.7)--cycle;
\draw[photon] (-1.8,1.5)--(-.7,.7);\draw[photon] (1.8,1.5)--(.7,.7);\draw[photon] (-1.8,-1.5)--(-.7,-.7);\draw[photon] (1.8,-1.5)--(.7,-.7);
\node[below] at (0,-1.9){light--light scattering};
\end{scope}
\end{tikzpicture}

 \caption{QED 的基本一圈结构：电子自能、真空极化、顶角修正，以及首次产生光子--光子散射的电子盒图。}
 \label{fig:qed-loops}
\end{figure}
